\documentclass[12pt]{ximera}
\title{Homework 3: Pairwise Comparisons and Fairness Criteria}
\begin{document}
\begin{abstract}
Sections 1.5--1.6: the Method of Pairwise Comparisons, Condorcet candidates, and an
election where correcting a clerical error in the winner's favor causes the winner to
lose --- a violation of the monotonicity criterion.
\end{abstract}
\maketitle

\section*{Sections 1.5--1.6: Pairwise Comparisons and Arrow's Fairness Criteria}

\begin{problem}
Table 1 shows the preference schedule for Best Sandwich Shop of Central Story County.
The candidates are A, B, C, and D.

\begin{center}
\begin{tabular}{c|ccccc}
Preference & 16 & 12 & 8 & 4 & 2 \\\hline
1st & C & A & D & B & A \\
2nd & B & B & A & D & D \\
3rd & D & C & C & A & B \\
4th & A & D & B & C & C
\end{tabular}

\smallskip
Table 1. Preference schedule for Best Sandwich Shop.
\end{center}

\begin{prompt}
\textbf{(a)} How many pairwise comparisons must be made for this election?

$\answer{6}$

\begin{hint}
Every unordered pair of candidates is compared exactly once.
\end{hint}

\begin{solution}
With $4$ candidates, the number of pairs is
\[ C(4,2)=\frac{4\cdot 3}{2}=6. \]
\end{solution}
\end{prompt}

\begin{prompt}
\textbf{(b)} Write out the pairwise comparisons and total the points for each
candidate. In each matchup, enter the number of voters preferring the first-named
candidate.

A vs.\ B: A gets $\answer{22}$, B gets $\answer{20}$ \\
A vs.\ C: A gets $\answer{26}$, C gets $\answer{16}$ \\
A vs.\ D: A gets $\answer{14}$, D gets $\answer{28}$ \\
B vs.\ C: B gets $\answer{18}$, C gets $\answer{24}$ \\
B vs.\ D: B gets $\answer{32}$, D gets $\answer{10}$ \\
C vs.\ D: C gets $\answer{28}$, D gets $\answer{14}$

\smallskip
Total points --- A: $\answer{2}$ \quad B: $\answer{1}$ \quad C: $\answer{2}$ \quad D: $\answer{1}$

\begin{hint}
For each matchup, go column by column and ask which of the two candidates appears
higher on that column's ballot; award that column's voters to them. The winner of each
matchup gets $1$ point.
\end{hint}

\begin{solution}
For each pair, count the voters who rank the first candidate above the second. There
are $42$ voters, so the two numbers in each matchup add to $42$.
\[
\begin{array}{lll}
\text{A vs.\ B} & 22\text{--}20 & \text{A wins}\\
\text{A vs.\ C} & 26\text{--}16 & \text{A wins}\\
\text{A vs.\ D} & 14\text{--}28 & \text{D wins}\\
\text{B vs.\ C} & 18\text{--}24 & \text{C wins}\\
\text{B vs.\ D} & 32\text{--}10 & \text{B wins}\\
\text{C vs.\ D} & 28\text{--}14 & \text{C wins}
\end{array}
\]
Awarding $1$ point per matchup win: A $= 2$, C $= 2$, B $= 1$, D $= 1$. The points
total $6$, one per comparison. \checkmark
\end{solution}
\end{prompt}

\begin{prompt}
\textbf{(c)} Give the complete ranking of this election using the pairwise comparisons
method. In the case of a tie in points between two candidates, refer to their
head-to-head comparison from part (b) to break the tie.

1st: $\answer{A}$ \quad 2nd: $\answer{C}$ \quad 3rd: $\answer{B}$ \quad 4th: $\answer{D}$

\begin{solution}
The point totals leave two ties: A and C at $2$ points, and B and D at $1$ point.

Break each with the head-to-head result:
\begin{itemize}
\item A vs.\ C was $26$--$16$ for A, so A ranks above C.
\item B vs.\ D was $32$--$10$ for B, so B ranks above D.
\end{itemize}
The complete ranking is therefore A, C, B, D.

Note that D beat the eventual winner A head-to-head ($28$--$14$), yet finishes last.
\end{solution}
\end{prompt}
\end{problem}

\begin{problem}
Make up an example election where there exists a \emph{Condorcet candidate} who does
not win the election. Write out the preference schedule, identify which candidate is
the Condorcet candidate, state which vote tabulation method you are using, and
identify the winner under that method. What does this imply about the tabulation
method you chose?

\begin{freeResponse}
\end{freeResponse}

\begin{hint}
A Condorcet candidate is one who would beat every other candidate in a head-to-head
matchup. Try to build a schedule where one candidate is almost everyone's second
choice --- strong head-to-head, but few first-place votes.
\end{hint}

\begin{solution}
Many examples work. Here is a small one with three candidates and $9$ voters:

\begin{center}
\begin{tabular}{c|ccc}
Preference & 4 & 3 & 2 \\\hline
1st & A & C & B \\
2nd & B & B & C \\
3rd & C & A & A
\end{tabular}
\end{center}

\textbf{B is the Condorcet candidate.} Head-to-head:
\begin{itemize}
\item B vs.\ A: the $3$ and $2$ columns prefer B, so B wins $5$--$4$.
\item B vs.\ C: the $4$ and $2$ columns prefer B, so B wins $6$--$3$.
\end{itemize}
B beats everyone one-on-one.

\textbf{Under the Plurality method}, however, the first-place counts are A $4$,
C $3$, B $2$ --- so \textbf{A wins} and B finishes last, despite being preferred by a
majority over each rival individually.

\textbf{What this implies:} Plurality violates the \emph{Condorcet criterion}, which
says that if a Condorcet candidate exists, that candidate should win. The reason is
structural: B is almost everyone's second choice, and Plurality does not look past
first place.

This is not a defect unique to Plurality. Arrow's Impossibility Theorem says no ranked
voting method with three or more candidates can satisfy all of the fairness criteria
at once, so every method fails some criterion.
\end{solution}
\end{problem}

\begin{problem}
Table 2 gives a preference schedule for election to the Ketchup Advisory Board.

\begin{center}
\begin{tabular}{c|cccccc}
Preference & 15 & 11 & 8 & 5 & 4 & 2 \\\hline
1st & A & C & B & D & B & D \\
2nd & B & B & A & C & D & B \\
3rd & C & A & D & A & C & A \\
4th & D & D & C & B & A & C
\end{tabular}

\smallskip
Table 2. Preference schedule for the Ketchup Advisory Board.
\end{center}

\begin{prompt}
\textbf{(a)} Using Table 2, apply Plurality-with-Elimination to find the complete
ranking of the candidates.

\textbf{Round 1:} A: $\answer{15}$ \quad B: $\answer{12}$ \quad C: $\answer{11}$ \quad
D: $\answer{7}$ \qquad Eliminated: $\answer{D}$

\smallskip
\textbf{Round 2:} A: $\answer{15}$ \quad B: $\answer{14}$ \quad C: $\answer{16}$
\qquad Eliminated: $\answer{B}$

\smallskip
\textbf{Round 3:} A: $\answer{25}$ \quad C: $\answer{20}$

\smallskip
Winner: $\answer{A}$

\begin{solution}
There are $45$ voters.

\textbf{Round 1.} First-place votes: A $15$, B $8+4 = 12$, C $11$, D $5+2 = 7$.
D is eliminated.

\textbf{Round 2.} D's $7$ votes transfer: the $5$-voter column had C second, so those
go to C; the $2$-voter column had B second, so those go to B. Now
\[ \text{A } 15, \qquad \text{B } 12+2 = 14, \qquad \text{C } 11+5 = 16. \]
B is eliminated.

\textbf{Round 3.} B's $14$ votes transfer to the highest surviving candidate on each
of those ballots: the $8$-voter column has A next, the $4$-voter column has C next,
and the $2$-voter column (originally D, then B) has A next. So
\[ \text{A } = 15+8+2 = 25, \qquad \text{C } = 16+4 = 20. \]
\textbf{A wins}, and the complete ranking is A, C, B, D.
\end{solution}
\end{prompt}

\begin{prompt}
\textbf{(b)} Suppose a clerical error was made in Table 2, and five of the voters
actually preferred A above C. The corrected preference schedule is shown in Table 3
--- the only change is in the $5$-voter column, where A has moved up from 3rd to 2nd.

\begin{center}
\begin{tabular}{c|cccccc}
Preference & 15 & 11 & 8 & \textbf{5} & 4 & 2 \\\hline
1st & A & C & B & D & B & D \\
2nd & B & B & A & \textbf{A} & D & B \\
3rd & C & A & D & \textbf{C} & C & A \\
4th & D & D & C & B & A & C
\end{tabular}

\smallskip
Table 3. \emph{Corrected} preference schedule.
\end{center}

Apply Plurality-with-Elimination again to find the corrected ranking.

\textbf{Round 1:} A: $\answer{15}$ \quad B: $\answer{12}$ \quad C: $\answer{11}$ \quad
D: $\answer{7}$ \qquad Eliminated: $\answer{D}$

\smallskip
\textbf{Round 2:} A: $\answer{20}$ \quad B: $\answer{14}$ \quad C: $\answer{11}$
\qquad Eliminated: $\answer{C}$

\smallskip
\textbf{Round 3:} A: $\answer{20}$ \quad B: $\answer{25}$

\smallskip
Winner: $\answer{B}$

\begin{hint}
Round 1 is unchanged. The difference appears in round 2: when D is eliminated, where
do those $5$ votes go now?
\end{hint}

\begin{solution}
\textbf{Round 1} is identical: A $15$, B $12$, C $11$, D $7$. D is eliminated.

\textbf{Round 2} is where the correction matters. The $5$-voter column now reads
D, A, C, B, so when D is eliminated those $5$ votes go to \textbf{A} instead of C:
\[ \text{A } = 15+5 = 20, \qquad \text{B } = 12+2 = 14, \qquad \text{C } = 11. \]
Now \emph{C} has the fewest votes and is eliminated instead of B.

\textbf{Round 3.} C's $11$ votes (the column C, B, A, D) go to B:
\[ \text{A } = 20, \qquad \text{B } = 14+11 = 25. \]
\textbf{B wins}, and the corrected ranking is B, A, C, D.
\end{solution}
\end{prompt}

\begin{prompt}
\textbf{(c)} Which fairness criterion is violated by this election?

\begin{multipleChoice}
\choice{The majority criterion}
\choice{The Condorcet criterion}
\choice[correct]{The monotonicity criterion}
\choice{Independence of irrelevant alternatives}
\end{multipleChoice}

\begin{solution}
The \textbf{monotonicity criterion}, which says that if a candidate wins an election
and then some voters change their ballots to rank that candidate \emph{higher} (with
no other changes), that candidate should still win.

Here A won under Table 2. The only change in Table 3 is that five voters moved A
\emph{up}, from 3rd place to 2nd. Nothing else changed --- and A lost.

The mechanism is worth understanding: gaining those five votes in round 2 pushed A's
total from $15$ to $20$, which knocked C out of the race instead of B. But C was the
candidate A could beat in a final round; B was not. A's extra support changed
\emph{who A had to face}, and that was enough to cost A the election.

This is a known weakness of Plurality-with-Elimination (also called instant-runoff
voting).
\end{solution}
\end{prompt}
\end{problem}

\end{document}
